% !TeX root = ../main.tex
\section{研究区域、数据体系与预处理}\label{sec:data}

\subsection{研究区域与空间单元划分}\label{sec:study-area}
% 确定高位冰崩源区、泥石流沟道、下游影响区域，以及计算使用的子流域或网格单元。
\writingplaceholder

\subsection{多源数据及其时空特征}\label{sec:data-sources}
% 集中介绍 MET、RADM、SOIL、FLUX、DEM、冰川、历史灾点、土地覆盖以及新增遥感和天气预报数据。
\writingplaceholder

\subsection{数据质量控制与统一}\label{sec:quality-control}
% 包括缺失值、异常值、QC 筛选、时间对齐、坐标转换、空间重采样等。
\writingplaceholder

\subsection{灾害事件时间轴与对照样本构建}\label{sec:event-timeline}
% 明确灾害源区触发时刻、下游受灾时刻的区别，构建历史同期与非事件对照窗口。
\writingplaceholder

\subsection{模型基本假设与符号约定}\label{sec:assumptions-symbols}
% 只列真正影响数学模型的假设。通用符号集中说明，其余参数在首次使用时定义。
\writingplaceholder
